\documentclass{article}
\usepackage[utf8]{inputenc}
\usepackage{amsmath,amssymb}
\usepackage[most]{tcolorbox}
\usepackage{geometry}
\geometry{margin=2.5cm}

\newcommand{\step}[1]{\par\noindent\hangindent=3.5em\hangafter=1%
  \makebox[3.5em][l]{\textbf{#1.}}}
\newcommand{\steptag}[1]{\hfill{\small\textsf{[#1]}}}

\newtcolorbox{proofbox}[1]{
  colback=gray!5, colframe=gray!60,
  fonttitle=\bfseries, title={#1},
  breakable, enhanced,
  left=4pt, right=4pt, top=4pt, bottom=4pt,
  before skip=10pt, after skip=10pt
}

\begin{document}

\begin{proofbox}{Depth-d halo collapse: for an exact signed idempotent P w...}
\step{1} Depth-d halo collapse: for an exact signed idempotent P with 0 < delta(P) <= 1/4 and nonempty visible set W(P), any row index v with d = dist\_1(p\_v, conv\{p\_w : w in W\}), and any halo width a > 0, writing d\_j = dist\_1(p\_j, conv W), sigma\_a(v) = sum over \{j : d\_j > a*tau\} of max(P\_vj, 0) (tau = sqrt(delta)), sigma(v) = sum over \{j : d\_j > 0\} of max(P\_vj, 0), nu\_v the row negative mass, and C\_a(v) = sum over \{j : d\_j > a*tau, d\_j > d\} of max(P\_vj, 0)*(d\_j - d), one has d*(1 - sigma\_a(v)) <= (sigma(v) - sigma\_a(v))*a*tau + nu\_v*(2 + 4*delta) + C\_a(v). \steptag{validated} \\[6pt]
\step{1.1} Setup and bookkeeping: let C=conv\{p\_w:w in W\}, a\_j=P\_vj, a\_j\textasciicircum{}+=max(a\_j,0), a\_j\textasciicircum{}-=max(-a\_j,0), nu\_v=sum\_j a\_j\textasciicircum{}-, H=\{j:d\_j>a*tau\}, and B=\{j:d\_j<=a*tau\}. From def-signed-idempotent, row v of P\textasciicircum{}2=P gives p\_v=sum\_j a\_j p\_j=sum\_j a\_j\textasciicircum{}+ p\_j-sum\_j a\_j\textasciicircum{}- p\_j; by lem-mass-split, sum\_j a\_j\textasciicircum{}+=1+nu\_v. Also d,d\_j>=0, H subset \{j:d\_j>0\}, sigma\_a<=sigma, C\_a>=0, and for every x in C and every row k, ||x-p\_k||\_1<=2+4*delta by the row-geometry clause of def-signed-idempotent and convexity of C. \steptag{validated} \\[4pt]
\step{1.2} Trivial branch: if sigma\_a(v)>=1, then d*(1-sigma\_a(v))<=0, while (sigma(v)-sigma\_a(v))*a*tau, nu\_v*(2+4*delta), and C\_a(v) are all nonnegative by the setup facts and a>0, tau=sqrt(delta)>0, nu\_v>=0; hence the required inequality holds in this branch. \steptag{validated} \\[4pt]
\step{1.3} Residual construction in the nontrivial branch: if sigma\_a(v)<1 and m=1-sigma\_a(v), define q=(sum\_\{j in B\} a\_j\textasciicircum{}+ p\_j - sum\_j a\_j\textasciicircum{}- p\_j)/m. Then m>0 and p\_v=sum\_\{j in H\} a\_j\textasciicircum{}+ p\_j + m*q; the coefficients a\_j\textasciicircum{}+ on H together with m are nonnegative and sum to 1. \steptag{validated} \\[4pt]
\step{1.4} Convex-distance rearrangement in the nontrivial branch: under sigma\_a(v)<1 and q as above, (1-sigma\_a(v))*d <= C\_a(v) + (1-sigma\_a(v))*dist\_1(q,C). \steptag{validated} \\[6pt]
\step{1.4.1} Convexity step: in the branch sigma\_a(v)<1, node 1.3 gives p\_v=sum\_\{j in H\} a\_j\textasciicircum{}+ p\_j + m*q with m=1-sigma\_a(v), nonnegative weights, and total weight sigma\_a(v)+m=1. Since C is convex, the norm-distance x $|-\rangle$ dist\_1(x,C) is convex, so d=dist\_1(p\_v,C) <= sum\_\{j in H\} a\_j\textasciicircum{}+ d\_j + m*dist\_1(q,C). \steptag{validated} \\[4pt]
\step{1.4.2} Deeper-row correction: for H=\{j:d\_j>a*tau\}, split H into indices with d\_j<=d and d\_j>d. Then sum\_\{j in H\} a\_j\textasciicircum{}+ d\_j <= sigma\_a(v)*d + sum\_\{j in H,d\_j>d\} a\_j\textasciicircum{}+*(d\_j-d) = sigma\_a(v)*d + C\_a(v). Combining this with the convexity step and moving sigma\_a(v)*d to the left gives (1-sigma\_a(v))*d <= C\_a(v)+(1-sigma\_a(v))*dist\_1(q,C). \steptag{validated} \\[4pt]
\step{1.5} Residual pricing in the nontrivial branch: under sigma\_a(v)<1 and q as above, (1-sigma\_a(v))*dist\_1(q,C) <= (sigma(v)-sigma\_a(v))*a*tau + nu\_v*(2+4*delta). \steptag{validated} \\[6pt]
\step{1.5.1} Apply lem-residual-upper: in the branch sigma\_a(v)<1, take positives (b\_j,p\_j)=(a\_j\textasciicircum{}+,p\_j) for j in B=\{j:d\_j<=a*tau\}, negatives (c\_k,r\_k)=(a\_k\textasciicircum{}-,p\_k) for all k, m=sum\_\{j in B\}a\_j\textasciicircum{}+-sum\_k a\_k\textasciicircum{}-=1-sigma\_a(v), and q as in node 1.3. For every x in C and every k, ||x-p\_k||\_1<=2+4*delta by node 1.1, so with D\_k=2+4*delta, lem-residual-upper gives m*dist\_1(q,C) <= sum\_\{j in B\} a\_j\textasciicircum{}+ d\_j + nu\_v*(2+4*delta). \steptag{validated} \\[4pt]
\step{1.5.2} Price the non-halo positive rows: on B one has d\_j<=a*tau, while rows with d\_j=0 contribute zero. Since H=\{j:d\_j>a*tau\} and sigma(v)-sigma\_a(v)=sum\_\{j:0<d\_j<=a*tau\} a\_j\textasciicircum{}+, it follows that sum\_\{j in B\} a\_j\textasciicircum{}+ d\_j <= a*tau*(sigma(v)-sigma\_a(v)). Combining this with the lem-residual-upper inequality proves the residual-pricing claim. \steptag{validated} \\[4pt]
\step{1.6} Branch assembly: the trivial branch proves the root when sigma\_a(v)>=1; when sigma\_a(v)<1, adding the convex-distance rearrangement and residual-pricing inequalities gives d*(1-sigma\_a(v)) <= C\_a(v)+(sigma(v)-sigma\_a(v))*a*tau+nu\_v*(2+4*delta), which is exactly the root inequality. \steptag{validated} \\[4pt]
\end{proofbox}

\end{document}
